\documentclass[11pt]{article}

\usepackage[utf8]{inputenc}
\usepackage[T1]{fontenc}
\usepackage{lmodern}
\usepackage{geometry}
\usepackage{amsmath,amssymb,amsfonts,amsthm,mathtools}
\usepackage{bm}
\usepackage{enumitem}
\usepackage{microtype}
\usepackage{hyperref}
\usepackage{titlesec}
\usepackage{booktabs}
\usepackage{array}
\usepackage{setspace}
\usepackage{mathrsfs}
\usepackage{physics}

\geometry{margin=1in}
\hypersetup{colorlinks=true, linkcolor=black, urlcolor=blue, citecolor=black}
\setlist[enumerate]{topsep=0.4em,itemsep=0.35em}
\setlist[itemize]{topsep=0.3em,itemsep=0.25em}
\titleformat{\section}{\large\bfseries}{\thesection.}{0.5em}{}
\titleformat{\subsection}{\normalsize\bfseries}{\thesubsection.}{0.5em}{}
\setstretch{1.06}

\newcommand{\Aether}{\AE{}ther}
\newcommand{\Aflow}{\AE{}ther-flow}
\newcommand{\TheoryName}{\Aflow{} Interpretation of Relativity}
\newcommand{\TheoryProgram}{\Aether{} / \Aflow{} framework}

\newcommand{\CompletedMetric}{g^{\sharp}}

\newcommand{\Car}{\mathrm{car}}
\newcommand{\Cur}{\mathrm{cur}}
\newcommand{\Hom}{\mathrm{hom}}

\newcommand{\ForSpace}[1]{\mathcal I_{#1}^{\mathrm{for}}}
\newcommand{\PrimShell}{\mathcal S_{\mathrm{prim}}}

\newcommand{\Reduction}[1]{\mathcal R_{#1}}
\newcommand{\CurrentCarrier}[1]{\mathfrak C_{#1}^{\mathrm{cur}}}
\newcommand{\EqCarrier}[1]{\mathfrak C_{#1}^{\mathrm{eq}}}

\newcommand{\MetricOut}{\mathscr G_{\mathrm{out}}}
\newcommand{\MatterOut}{\mathscr T_{\mathrm{out}}}

\newcommand{\VisibleAffine}[1]{\widehat X_{#1}}
\newcommand{\DataSpace}[1]{\mathcal Y_{#1}^{\mathrm{obs}}}
\newcommand{\AmbientTarget}[1]{E_{#1}^{\mathrm{amb}}}

\newcommand{\EqnSpace}[1]{\mathcal U_{#1}^{\mathrm{eqn}}}
\newcommand{\EqnBody}[1]{\mathcal B_{#1}^{\mathrm{eqn}}}
\newcommand{\EqnData}[1]{\Lambda_{#1}^{\mathrm{eqn,dat}}}
\newcommand{\EqnShell}[1]{\Lambda_{#1}^{\mathrm{eqn,sh}}}
\newcommand{\EqnTarget}[1]{Q_{#1}^{\mathrm{eqn}}}
\newcommand{\EqnPair}[1]{\mathfrak b_{#1}^{\mathrm{eqn}}}
\newcommand{\EqnHidden}[1]{H_{#1}^{\mathrm{eqn}}}
\newcommand{\EqnZero}[1]{V_{#1}^{\mathrm{eqn},0}}

\newcommand{\ResSpace}[1]{\mathcal U_{#1}^{\mathrm{sup,res}}}
\newcommand{\ResBody}[1]{\mathcal B_{#1}^{\mathrm{sup,res}}}
\newcommand{\ResTarget}[1]{S_{#1}^{\mathrm{sup,res}}}
\newcommand{\ResPair}[1]{\mathfrak m_{#1}^{\mathrm{sup,res}}}
\newcommand{\ResSection}[1]{\Theta_{#1}^{\mathrm{sup,res}}}
\newcommand{\ResData}[1]{\Lambda_{#1}^{\mathrm{sup,dat}}}
\newcommand{\ResShellMap}[1]{\Lambda_{#1}^{\mathrm{sup,sh}}}
\newcommand{\ResShellSet}[1]{\mathcal Z_{#1}^{\mathrm{sup,res}}}
\newcommand{\SeedHidden}[1]{\mathcal H_{#1}^{\mathrm{sup}}}
\newcommand{\SeedHiddenBody}[1]{D_{#1}^{\mathrm{sup}}}
\newcommand{\HiddenChart}[1]{\Psi_{#1}^{\mathrm{hid}}}
\newcommand{\ZeroBody}[1]{\mathcal B_{#1}^{0,\mathrm{sup}}}

\newcommand{\SubSpace}[1]{\mathcal M_{#1}^{\mathrm{sub}}}
\newcommand{\SubBody}[1]{\mathcal B_{#1}^{\mathrm{sub}}}
\newcommand{\SubTarget}[1]{E_{#1}^{\mathrm{sub}}}
\newcommand{\SubPair}[1]{\mathfrak s_{#1}^{\mathrm{sub}}}
\newcommand{\SubShellSet}[1]{\mathcal Z_{#1}^{\mathrm{sub}}}

\newcommand{\LiftSpace}[1]{\mathcal L_{#1}^{\mathrm{lift}}}
\newcommand{\LiftBody}[1]{\mathcal B_{#1}^{\mathrm{lift}}}
\newcommand{\LiftProj}[1]{\Pi_{#1}^{\mathrm{lift}}}
\newcommand{\LiftSelect}[1]{\mathfrak s_{#1}^{\mathrm{lift}}}
\newcommand{\LiftSection}[1]{\Gamma_{#1}^{\mathrm{lift}}}
\newcommand{\LiftFunctional}[1]{\mathscr V_{#1}^{\mathrm{lift}}}
\newcommand{\LiftData}[1]{\Lambda_{#1}^{\mathrm{lift,dat}}}
\newcommand{\LiftShellMap}[1]{\Lambda_{#1}^{\mathrm{lift,sh}}}
\newcommand{\LiftSym}[1]{\mathbb G_{#1}^{\mathrm{lift}}}
\newcommand{\LiftTime}[1]{\vartheta_{#1}^{\mathrm{lift}}}
\newcommand{\LiftShellSet}[1]{\mathcal Z_{#1}^{\mathrm{lift}}}
\newcommand{\LiftChart}[1]{\Xi_{#1}^{\mathrm{lift}}}
\newcommand{\LiftSupport}[1]{\Theta_{#1}^{\mathrm{lift,sup}}}
\newcommand{\LiftSupportShell}[1]{\mathcal Z_{#1}^{\mathrm{lift,sup}}}
\newcommand{\LiftZeroCore}[1]{\mathcal Z_{#1}^{0,\mathrm{lift}}}
\newcommand{\LiftEqChart}[1]{\Psi_{#1}^{\mathrm{lift,sup}}}
\newcommand{\LiftZeroChart}[1]{\Psi_{#1}^{0,\mathrm{lift}}}
\newcommand{\LiftSplit}[1]{\Phi_{#1}^{\mathrm{lift}}}
\newcommand{\LiftCharge}[1]{\mathcal Q_{#1}^{\mathrm{lift}}}
\newcommand{\LiftResidual}[1]{\mathcal R_{#1}^{\mathrm{lift}}}
\newcommand{\LiftEmbed}[1]{\zeta_{#1}^{\mathrm{lift}}}
\newcommand{\LiftKernel}[1]{K_{#1}^{0,\mathrm{lift}}}
\newcommand{\LiftKernelCh}[1]{\widetilde K_{#1}^{0,\mathrm{lift}}}
\newcommand{\LiftStateComp}[1]{P_{#1}^{\mathrm{lift,co}}}

\newcommand{\PrecSpace}[1]{\mathcal P_{#1}^{\mathrm{prec}}}
\newcommand{\PrecBody}[1]{\mathcal B_{#1}^{\mathrm{prec}}}
\newcommand{\PrecProj}[1]{\Pi_{#1}^{\mathrm{prec}}}
\newcommand{\PrecLiftMin}[1]{\mathfrak f_{#1}^{\mathrm{prec}}}
\newcommand{\PrecSubMin}[1]{\mathfrak m_{#1}^{\mathrm{prec}}}
\newcommand{\PrecSection}[1]{\Gamma_{#1}^{\mathrm{prec}}}
\newcommand{\PrecFunctional}[1]{\mathscr W_{#1}^{\mathrm{prec}}}
\newcommand{\PrecData}[1]{\Lambda_{#1}^{\mathrm{prec,dat}}}
\newcommand{\PrecShellMap}[1]{\Lambda_{#1}^{\mathrm{prec,sh}}}
\newcommand{\PrecSym}[1]{\mathbb G_{#1}^{\mathrm{prec}}}
\newcommand{\PrecTime}[1]{\vartheta_{#1}^{\mathrm{prec}}}
\newcommand{\PrecShellSet}[1]{\mathcal Z_{#1}^{\mathrm{prec}}}
\newcommand{\PrecChart}[1]{\Xi_{#1}^{\mathrm{prec}}}
\newcommand{\PrecSupport}[1]{\Theta_{#1}^{\mathrm{prec,sup}}}
\newcommand{\PrecSupportShell}[1]{\mathcal Z_{#1}^{\mathrm{prec,sup}}}
\newcommand{\PrecZeroCore}[1]{\mathcal Z_{#1}^{0,\mathrm{prec}}}
\newcommand{\PrecEqChart}[1]{\Psi_{#1}^{\mathrm{prec,sup}}}
\newcommand{\PrecZeroChart}[1]{\Psi_{#1}^{0,\mathrm{prec}}}
\newcommand{\PrecSplit}[1]{\Phi_{#1}^{\mathrm{prec}}}
\newcommand{\PrecCharge}[1]{\mathcal Q_{#1}^{\mathrm{prec}}}
\newcommand{\PrecResidual}[1]{\mathcal R_{#1}^{\mathrm{prec}}}
\newcommand{\PrecEmbed}[1]{\zeta_{#1}^{\mathrm{prec}}}
\newcommand{\PrecKernel}[1]{K_{#1}^{0,\mathrm{prec}}}
\newcommand{\PrecKernelCh}[1]{\widetilde K_{#1}^{0,\mathrm{prec}}}
\newcommand{\PrecStateComp}[1]{P_{#1}^{\mathrm{prec,co}}}

\newtheorem{definition}{Definition}
\newtheorem{lemma}{Lemma}
\newtheorem{proposition}{Proposition}
\newtheorem{remark}{Remark}
\newtheorem{corollary}{Corollary}
\newtheorem{theorem}{Theorem}

\title{\textbf{The \TheoryName{}}\\[0.4em]
\large Primitive-Reservoir Observer-Reduction\\
\large Shell-Generating Substrate Lift Dynamics\\
\large Necessity / Uniqueness from Lift-Generating Substrate Precursor Dynamics}
\author{}
\date{}

\begin{document}

\maketitle

\begin{abstract}
The current benchmark-facing frontier on the active primitive-reservoir observer-reduction line is no longer the split selector-shell core itself. The shell-restriction forcing theorem has already spent that move: the split zero-hidden selector-shell core is now a canonical and unique output of primitive substrate shell restriction. The next honest benchmark-facing burden is therefore one level lower. Its first admissible target is the shell-generating substrate lift dynamics that support the shell-restriction layer.

The recorded lift-origin theorem already showed that this lift package can be produced canonically from a deeper lift-generating substrate precursor class. The present note sharpens that existence statement to the benchmark-facing forcing verdict that the routing layer now requires. Starting from lift-generating substrate precursor dynamics, one recovers canonically the lifted restriction section, the exact lifted shell, the fiber-minimizing substrate lift, the lifted shell/support charts, the lifted support zero core, the zero-core / hidden-seed split, the hidden charge and hidden recorder, the observer-compatible exact zero-response realization, the fixed-data solvability statement, and the coercive control on lifted data-invisible directions. Within that precursor class the shell-generating lift package is not merely available. It is necessary and unique up to the obvious precursor-faithful identification preserving the same lifted fibers, the same charted support shell, the same zero-hidden split and hidden quadratic form, and the same benchmark shell data.

The benchmark boundary remains fixed throughout. The one operative metric is still exactly \(\CompletedMetric\), the ordinary matter route is unchanged, there is no second operative metric, the low-energy matter law is not enlarged, and stronger first-principles promotion language remains paused. The honest burden therefore moves below lift scope to the precursor layer itself.
\end{abstract}

\section{Introduction}

The active derivational program of \TheoryName{} has now reached the point where another shell-restriction relay would be benchmark neutral. The current shell-restriction forcing theorem already removed the split zero-hidden selector-shell core as the live stopping point. Once that theorem is on record, the next honest benchmark-facing move must go one level lower and address the shell-generating substrate lift dynamics that support the shell-restriction layer.

There are two disciplined ways to do that. One can derive the lift package from still more primitive explicit substrate structure, or one can prove that the lift package is already forced inside a sharper deeper class. The recorded precursor-origin theorem supplied the first ingredient needed for the second route: it exhibited the entire lift package as a projected exact-shell transport of a deeper lift-generating substrate precursor class. What it did not yet state sharply enough for the current routing burden is the corresponding necessity / uniqueness verdict at lift scope.

The present note records exactly that stronger claim. It does not change the benchmark exact-closure package. It does not introduce a second operative metric or a new low-energy matter law. It only proves that, once the precursor package is fixed, no independent benchmark-facing freedom remains at shell-generating lift scope.

\paragraph{Framework and claim boundary.}
\TheoryProgram{} denotes the broader \Aether{} / \Aflow{} framework built on the claims that the \Aether{} is the underlying four-dimensional substrate of reality, the \Aflow{} is its intrinsic ordered motion, observed three-dimensional space is the local experiential slice of that deeper substrate, S-time is the experienced order of change arising from matter, light, and the \Aflow{}, and observed expansion is the three-dimensional appearance of deeper four-dimensional ordered motion. Within that framework, \TheoryName{} denotes the exact-closure theory in which the effective gravitational dynamics are adopted to be exactly Einsteinian with universal matter coupling. In the active sequence, \emph{adoption} means use of that established relativistic dynamics without claiming substrate derivation, while \emph{derivation} is reserved for a first-principles recovery from explicit substrate variables.

The present note stays strictly inside that boundary. It does not claim that final substrate microphysics has already been identified. Its narrower theorem is only that, on the active primitive-reservoir observer-reduction line, the shell-generating substrate lift dynamics are themselves a forced and unique output of the recorded precursor class that sits immediately below them.

\section{Fixed Primitive Continuation Data}

\begin{definition}[Recorded primitive continuation data]
\label{def:fixed}
Fix the following continuation data from the active primitive-reservoir observer-reduction line.
\begin{enumerate}
    \item For each sector label \(\sigma\in\{\Car,\Cur,\Hom\}\), a primitive forcing shell
    \[
    \ForSpace{\sigma},
    \]
    a visible reduction map
    \[
    \Reduction{\sigma}:\ForSpace{\sigma}\to X_{\sigma},
    \]
    and shell-level current and equilibrium carriers
    \[
    \CurrentCarrier{\sigma}:\ForSpace{\sigma}\to Z_{\sigma}^{\mathrm{cur}},
    \qquad
    \EqCarrier{\sigma}:\ForSpace{\sigma}\to Z_{\sigma}^{\mathrm{eq}}.
    \]
    \item The product shell
    \[
    \PrimShell
    \equiv
    \ForSpace{\Car}\times \ForSpace{\Cur}\times \ForSpace{\Hom}.
    \]
    \item The recorded metric and ordinary-matter outputs
    \[
    \MetricOut:\PrimShell\to\{\text{observer metrics}\},
    \qquad
    \MatterOut:\PrimShell\times\Psi\to\{\text{observer stress tensors}\},
    \]
    satisfying
    \begin{equation}
    \MetricOut(\mathbf w)=\CompletedMetric,
    \qquad
    \MatterOut(\mathbf w,\psi)
    =
    T_{\mu\nu}^{\mathrm{matter}}[\CompletedMetric,\psi]
    \label{eq:fixedoutputs}
    \end{equation}
    for every \(\mathbf w\in\PrimShell\) and every matter configuration \(\psi\in\Psi\).
\end{enumerate}
\end{definition}

\begin{remark}
Definition \ref{def:fixed} is not reopened here. The primitive continuation data, the one operative metric, and the ordinary matter route remain fixed continuation data. The present note only removes the shell-generating lift package itself as the live unexplained stopping point.
\end{remark}

\section{Shell-Generating Substrate Lift Dynamics}

\begin{definition}[Shell-generating substrate lift dynamics]
\label{def:lift}
Fix \(\sigma\in\{\Car,\Cur,\Hom\}\). A \emph{shell-generating substrate lift dynamics package} for sector \(\sigma\) consists of the following data.
\begin{enumerate}
    \item A finite-dimensional Euclidean lifted state space
    \[
    \LiftSpace{\sigma},
    \]
    a compact convex lifted body
    \[
    \LiftBody{\sigma}\subset\LiftSpace{\sigma}
    \]
    with nonempty interior, an affine surjection
    \[
    \LiftProj{\sigma}:\LiftSpace{\sigma}\to\SubSpace{\sigma},
    \]
    and the projected substrate body
    \[
    \SubBody{\sigma}
    \equiv
    \LiftProj{\sigma}\bigl(\LiftBody{\sigma}\bigr)
    \subset \SubSpace{\sigma},
    \]
    assumed compact convex with nonempty interior. There is a smooth selection
    \[
    \LiftSelect{\sigma}:\SubBody{\sigma}\to\LiftBody{\sigma}
    \]
    satisfying
    \[
    \LiftProj{\sigma}\circ\LiftSelect{\sigma}
    =
    \mathrm{id}_{\SubBody{\sigma}}.
    \]
    \item The same finite-dimensional Euclidean substrate restriction target
    \[
    \SubTarget{\sigma},
    \]
    the same positive-definite symmetric substrate pairing
    \[
    \SubPair{\sigma}:
    \SubTarget{\sigma}\times\SubTarget{\sigma}\to\mathbb R,
    \]
    and a smooth lifted restriction section
    \[
    \LiftSection{\sigma}:\LiftSpace{\sigma}\to\SubTarget{\sigma}.
    \]
    Define the lifted restriction functional
    \begin{equation}
    \LiftFunctional{\sigma}(\ell)
    \equiv
    \SubPair{\sigma}\bigl(\LiftSection{\sigma}(\ell),\LiftSection{\sigma}(\ell)\bigr).
    \label{eq:liftfunctional}
    \end{equation}
    Define the exact lifted shell
    \begin{equation}
    \LiftShellSet{\sigma}
    \equiv
    \bigl(\LiftSection{\sigma}\bigr)^{-1}(0)\cap\LiftBody{\sigma}.
    \label{eq:liftshell}
    \end{equation}
    For each \(m\in\SubBody{\sigma}\), assume \(\LiftSelect{\sigma}(m)\) is the unique minimizer of \(\LiftFunctional{\sigma}\) on the fiber
    \[
    \bigl(\LiftProj{\sigma}\bigr)^{-1}(m)\cap\LiftBody{\sigma},
    \]
    and that there exists \(\mu_{\sigma}^{\mathrm{fib}}>0\) such that
    \begin{equation}
    \LiftFunctional{\sigma}(\ell)
    \ge
    \LiftFunctional{\sigma}\bigl(\LiftSelect{\sigma}(m)\bigr)
    +
    \mu_{\sigma}^{\mathrm{fib}}\,
    \operatorname{dist}\bigl(\ell,\LiftSelect{\sigma}(m)\bigr)^{2}
    \label{eq:fibergap}
    \end{equation}
    for every \(m\in\SubBody{\sigma}\) and every \(\ell\in\bigl(\LiftProj{\sigma}\bigr)^{-1}(m)\cap\LiftBody{\sigma}\). Assume there exists \(\mu_{\sigma}^{\mathrm{sub}}>0\) such that
    \begin{equation}
    \LiftFunctional{\sigma}\bigl(\LiftSelect{\sigma}(m)\bigr)
    \ge
    \mu_{\sigma}^{\mathrm{sub}}\,
    \operatorname{dist}\bigl(m,\LiftProj{\sigma}(\LiftShellSet{\sigma})\bigr)^{2}
    \qquad
    \text{for all }m\in\SubBody{\sigma}.
    \label{eq:projectedgap}
    \end{equation}
    \item Affine lifted observer-data and shell-response readouts
    \[
    \LiftData{\sigma}:\LiftSpace{\sigma}\to\DataSpace{\sigma},
    \qquad
    \LiftShellMap{\sigma}:\LiftSpace{\sigma}\to\AmbientTarget{\sigma},
    \]
    a finite-dimensional Euclidean support target
    \[
    \ResTarget{\sigma},
    \]
    a positive-definite symmetric support pairing
    \[
    \ResPair{\sigma}:
    \ResTarget{\sigma}\times\ResTarget{\sigma}\to\mathbb R,
    \]
    a finite-dimensional real hidden target
    \[
    \EqnTarget{\sigma},
    \]
    and a positive-definite symmetric hidden pairing
    \[
    \EqnPair{\sigma}:
    \EqnTarget{\sigma}\times\EqnTarget{\sigma}\to\mathbb R.
    \]
    There is also a compact symmetry group
    \[
    \LiftSym{\sigma}
    \]
    and an involution
    \[
    \LiftTime{\sigma},
    \]
    acting orthogonally on \(\LiftSpace{\sigma}\), \(\SubSpace{\sigma}\), \(\DataSpace{\sigma}\), \(\AmbientTarget{\sigma}\), and \(\ResTarget{\sigma}\), and by \(\EqnPair{\sigma}\)-isometries on \(\EqnTarget{\sigma}\), preserving \(\LiftBody{\sigma}\), \(\LiftProj{\sigma}\), \(\LiftSelect{\sigma}\), \(\LiftSection{\sigma}\), \(\LiftData{\sigma}\), and \(\LiftShellMap{\sigma}\).
    \item The restriction of \(\LiftProj{\sigma}\) to \(\LiftShellSet{\sigma}\) is a diffeomorphism onto its image
    \[
    \SubBody{\sigma}
    \supset
    \SubShellSet{\sigma}
    \equiv
    \LiftProj{\sigma}\bigl(\LiftShellSet{\sigma}\bigr).
    \]
    There is a Euclidean lifted shell chart
    \[
    \LiftChart{\sigma}:\LiftShellSet{\sigma}\to\ResSpace{\sigma}
    \]
    whose shell-body image
    \[
    \ResBody{\sigma}
    \equiv
    \LiftChart{\sigma}\bigl(\LiftShellSet{\sigma}\bigr)
    \]
    is compact convex with nonempty interior. There exist affine maps
    \[
    \ResData{\sigma}:\ResSpace{\sigma}\to\DataSpace{\sigma},
    \qquad
    \ResShellMap{\sigma}:\ResSpace{\sigma}\to\AmbientTarget{\sigma},
    \]
    and a smooth support section
    \[
    \ResSection{\sigma}:\ResSpace{\sigma}\to\ResTarget{\sigma}
    \]
    such that
    \begin{equation}
    \ResData{\sigma}\circ\LiftChart{\sigma}
    =
    \LiftData{\sigma}\big|_{\LiftShellSet{\sigma}},
    \qquad
    \ResShellMap{\sigma}\circ\LiftChart{\sigma}
    =
    \LiftShellMap{\sigma}\big|_{\LiftShellSet{\sigma}}.
    \label{eq:liftchartreadouts}
    \end{equation}
    There is a smooth lifted shell-internal support recorder
    \[
    \LiftSupport{\sigma}:\LiftShellSet{\sigma}\to\ResTarget{\sigma}
    \]
    satisfying
    \begin{equation}
    \ResSection{\sigma}\circ\LiftChart{\sigma}
    =
    \LiftSupport{\sigma}.
    \label{eq:liftsupporttransport}
    \end{equation}
    Define the exact lifted support shell
    \begin{equation}
    \LiftSupportShell{\sigma}
    \equiv
    \bigl(\LiftSupport{\sigma}\bigr)^{-1}(0)\cap\LiftShellSet{\sigma},
    \label{eq:liftsupportshell}
    \end{equation}
    and assume
    \begin{equation}
    \ResShellSet{\sigma}
    \equiv
    \bigl(\ResSection{\sigma}\bigr)^{-1}(0)\cap\ResBody{\sigma}
    =
    \LiftChart{\sigma}\bigl(\LiftSupportShell{\sigma}\bigr).
    \label{eq:lifttransportedsupportshell}
    \end{equation}
    Assume there exists \(\nu_{\sigma}^{\mathrm{sup}}>0\) such that
    \begin{equation}
    \ResPair{\sigma}\bigl(\ResSection{\sigma}(u),\ResSection{\sigma}(u)\bigr)
    \ge
    \nu_{\sigma}^{\mathrm{sup}}\,
    \operatorname{dist}\bigl(u,\ResShellSet{\sigma}\bigr)^{2}
    \qquad
    \text{for all }u\in\ResBody{\sigma}.
    \label{eq:liftnormal}
    \end{equation}
    The transported symmetry and time-reversal actions act linearly on \(\ResSpace{\sigma}\), preserving \(\ResBody{\sigma}\), \(\ResShellSet{\sigma}\), \(\ResSection{\sigma}\), \(\ResData{\sigma}\), and \(\ResShellMap{\sigma}\).
    \item A closed embedded lifted support zero core
    \[
    \LiftZeroCore{\sigma}\subset\LiftSupportShell{\sigma},
    \]
    a finite-dimensional Euclidean hidden-seed space
    \[
    \SeedHidden{\sigma},
    \]
    and a compact convex hidden-seed body
    \[
    \SeedHiddenBody{\sigma}\subset\SeedHidden{\sigma}
    \]
    with
    \[
    0\in\operatorname{int}\SeedHiddenBody{\sigma}.
    \]
    There is a Euclidean lifted support-shell chart
    \[
    \LiftEqChart{\sigma}:\LiftSupportShell{\sigma}\to\EqnSpace{\sigma},
    \]
    an affine chart on the lifted support zero core
    \[
    \LiftZeroChart{\sigma}:\LiftZeroCore{\sigma}\to\EqnZero{\sigma},
    \]
    and a linear isometric embedding
    \[
    \HiddenChart{\sigma}:\SeedHidden{\sigma}\to\EqnHidden{\sigma}
    \]
    such that
    \[
    \EqnSpace{\sigma}
    =
    \EqnZero{\sigma}\oplus\EqnHidden{\sigma}
    \]
    is an orthogonal direct sum, the set
    \[
    \ZeroBody{\sigma}
    \equiv
    \LiftZeroChart{\sigma}\bigl(\LiftZeroCore{\sigma}\bigr)
    \]
    is compact convex with nonempty interior in \(\EqnZero{\sigma}\), the shell image
    \[
    \EqnBody{\sigma}
    \equiv
    \LiftEqChart{\sigma}\bigl(\LiftSupportShell{\sigma}\bigr)
    \]
    has nonempty interior in \(\EqnSpace{\sigma}\), and there is a diffeomorphism
    \begin{equation}
    \LiftSplit{\sigma}:
    \LiftZeroCore{\sigma}\times\SeedHiddenBody{\sigma}
    \xrightarrow{\cong}
    \LiftSupportShell{\sigma}
    \label{eq:liftsplit}
    \end{equation}
    satisfying
    \begin{equation}
    \LiftEqChart{\sigma}\bigl(\LiftSplit{\sigma}(z,\eta)\bigr)
    =
    \LiftZeroChart{\sigma}(z)+\HiddenChart{\sigma}(\eta).
    \label{eq:liftchartsplit}
    \end{equation}
    The lifted shell readouts become affine in the lifted support-shell chart: there exist affine maps
    \[
    \EqnData{\sigma}:\EqnSpace{\sigma}\to\DataSpace{\sigma},
    \qquad
    \EqnShell{\sigma}:\EqnSpace{\sigma}\to\AmbientTarget{\sigma},
    \]
    such that
    \begin{equation}
    \EqnData{\sigma}\circ\LiftEqChart{\sigma}
    =
    \LiftData{\sigma}\big|_{\LiftSupportShell{\sigma}},
    \qquad
    \EqnShell{\sigma}\circ\LiftEqChart{\sigma}
    =
    \LiftShellMap{\sigma}\big|_{\LiftSupportShell{\sigma}}.
    \label{eq:lifteqnchartreadouts}
    \end{equation}
    The transported symmetry and time-reversal actions act linearly on \(\EqnSpace{\sigma}\), preserving \(\EqnZero{\sigma}\), \(\EqnHidden{\sigma}\), and \(\EqnBody{\sigma}\).
    \item A linear hidden charge map
    \[
    \LiftCharge{\sigma}:\SeedHidden{\sigma}\to\EqnTarget{\sigma}
    \]
    such that there exists \(\lambda_{\sigma}^{\mathrm{eqn}}>0\) with
    \begin{equation}
    \EqnPair{\sigma}\bigl(\LiftCharge{\sigma}\eta,\LiftCharge{\sigma}\eta\bigr)
    \ge
    \lambda_{\sigma}^{\mathrm{eqn}}\,
    \|\eta\|^{2}
    \qquad
    \text{for all }\eta\in\SeedHidden{\sigma}.
    \label{eq:liftchargegap}
    \end{equation}
    Define the lifted support-shell hidden recorder by
    \begin{equation}
    \LiftResidual{\sigma}\bigl(\LiftSplit{\sigma}(z,\eta)\bigr)
    \equiv
    \LiftCharge{\sigma}(\eta).
    \label{eq:liftresidual}
    \end{equation}
    \item A smooth embedding
    \[
    \jmath_{\sigma}:X_{\sigma}\hookrightarrow\VisibleAffine{\sigma}
    \]
    and a smooth observer-compatible exact lifted zero-response realization
    \[
    \LiftEmbed{\sigma}:\ForSpace{\sigma}\hookrightarrow\LiftZeroCore{\sigma}
    \]
    whose image is exactly
    \[
    \LiftZeroCore{\sigma}
    \cap
    \bigl(\LiftShellMap{\sigma}\bigr)^{-1}(0),
    \]
    and such that
    \begin{equation}
    \LiftData{\sigma}\bigl(\LiftEmbed{\sigma}(w)\bigr)
    =
    \bigl(
    \jmath_{\sigma}(\Reduction{\sigma}(w)),
    \CurrentCarrier{\sigma}(w),
    \EqCarrier{\sigma}(w)
    \bigr)
    \label{eq:liftcompat}
    \end{equation}
    for every \(w\in\ForSpace{\sigma}\).
    \item Fixed-data solvability on the lifted support zero core:
    \begin{equation}
    \LiftData{\sigma}\bigl(\LiftZeroCore{\sigma}\bigr)
    =
    \LiftData{\sigma}\Bigl(
    \LiftZeroCore{\sigma}\cap
    \bigl(\LiftShellMap{\sigma}\bigr)^{-1}(0)
    \Bigr).
    \label{eq:liftsolvability}
    \end{equation}
    \item The inert lifted support-zero-core kernel
    \[
    \LiftKernel{\sigma}
    \equiv
    \ker\bigl(D\LiftData{\sigma}\big|_{\LiftZeroCore{\sigma}}\bigr)
    \cap
    \ker\bigl(D\LiftShellMap{\sigma}\big|_{\LiftZeroCore{\sigma}}\bigr)
    \]
    and the orthogonal projector
    \[
    \LiftStateComp{\sigma}:
    \EqnZero{\sigma}\to
    \bigl(\LiftKernelCh{\sigma}\bigr)^{\perp}\cap\EqnZero{\sigma},
    \qquad
    \LiftKernelCh{\sigma}
    \equiv
    D\LiftZeroChart{\sigma}\bigl(\LiftKernel{\sigma}\bigr).
    \]
    There exists \(\kappa_{\sigma}^{\mathrm{eqn}}>0\) such that for every
    \[
    w\in\ForSpace{\sigma},
    \qquad
    \delta \ell\in T_{\LiftEmbed{\sigma}(w)}\LiftSupportShell{\sigma}
    \]
    satisfying
    \[
    D\LiftData{\sigma}\,\delta \ell=0,
    \]
    if
    \[
    D\LiftEqChart{\sigma}(\delta \ell)=\delta z+\delta\eta,
    \qquad
    \delta z\in\EqnZero{\sigma},
    \quad
    \delta\eta\in\EqnHidden{\sigma},
    \]
    then
    \begin{equation}
    \|D\LiftShellMap{\sigma}\,\delta \ell\|^{2}
    +
    \EqnPair{\sigma}\bigl(
    \LiftCharge{\sigma}\bigl((\HiddenChart{\sigma})^{-1}\delta\eta\bigr),
    \LiftCharge{\sigma}\bigl((\HiddenChart{\sigma})^{-1}\delta\eta\bigr)
    \bigr)
    \ge
    \kappa_{\sigma}^{\mathrm{eqn}}
    \left(
    \|\LiftStateComp{\sigma}(\delta z)\|^{2}
    +
    \|\delta\eta\|^{2}
    \right).
    \label{eq:liftcoercive}
    \end{equation}
\end{enumerate}
\end{definition}

\begin{remark}
Definition \ref{def:lift} is the precise package that now has to be explained one layer deeper. The benchmark claim boundary says that such an explanation may sharpen the origin-side derivational line, but it may not change the one operative metric, enlarge the low-energy matter law, or blur the distinction between exact relativistic adoption and first-principles derivation.
\end{remark}

\section{Lift-Generating Substrate Precursor Dynamics}

\begin{definition}[Lift-generating substrate precursor dynamics]
\label{def:prec}
Fix \(\sigma\in\{\Car,\Cur,\Hom\}\). A \emph{lift-generating substrate precursor dynamics package} for sector \(\sigma\) consists of the following data.
\begin{enumerate}
    \item A finite-dimensional Euclidean precursor state space
    \[
    \PrecSpace{\sigma},
    \]
    a compact convex precursor body
    \[
    \PrecBody{\sigma}\subset\PrecSpace{\sigma}
    \]
    with nonempty interior, an affine surjection
    \[
    \PrecProj{\sigma}:\PrecSpace{\sigma}\to\LiftSpace{\sigma},
    \]
    and the projected lift body
    \[
    \LiftBody{\sigma}
    \equiv
    \PrecProj{\sigma}\bigl(\PrecBody{\sigma}\bigr)
    \subset\LiftSpace{\sigma},
    \]
    assumed compact convex with nonempty interior. There is an affine surjection
    \[
    \LiftProj{\sigma}:\LiftSpace{\sigma}\to\SubSpace{\sigma},
    \]
    and the projected substrate body
    \[
    \SubBody{\sigma}
    \equiv
    \LiftProj{\sigma}\bigl(\LiftBody{\sigma}\bigr)
    \subset\SubSpace{\sigma},
    \]
    also assumed compact convex with nonempty interior.

    There is a smooth precursor lift-fiber minimizer
    \[
    \PrecLiftMin{\sigma}:\LiftBody{\sigma}\to\PrecBody{\sigma}
    \]
    satisfying
    \[
    \PrecProj{\sigma}\circ\PrecLiftMin{\sigma}
    =
    \mathrm{id}_{\LiftBody{\sigma}},
    \]
    and a smooth precursor substrate-fiber minimizer
    \[
    \PrecSubMin{\sigma}:\SubBody{\sigma}\to\PrecBody{\sigma}
    \]
    satisfying
    \[
    \LiftProj{\sigma}\circ\PrecProj{\sigma}\circ\PrecSubMin{\sigma}
    =
    \mathrm{id}_{\SubBody{\sigma}}.
    \]
    Define
    \begin{equation}
    \LiftSelect{\sigma}
    \equiv
    \PrecProj{\sigma}\circ\PrecSubMin{\sigma}.
    \label{eq:inducedliftselect}
    \end{equation}
    Assume
    \begin{equation}
    \PrecLiftMin{\sigma}\bigl(\LiftSelect{\sigma}(m)\bigr)
    =
    \PrecSubMin{\sigma}(m)
    \qquad
    \text{for every }m\in\SubBody{\sigma}.
    \label{eq:compatmin}
    \end{equation}
    \item The same finite-dimensional Euclidean substrate restriction target
    \[
    \SubTarget{\sigma},
    \]
    the same positive-definite symmetric substrate pairing
    \[
    \SubPair{\sigma}:
    \SubTarget{\sigma}\times\SubTarget{\sigma}\to\mathbb R,
    \]
    and a smooth precursor restriction section
    \[
    \PrecSection{\sigma}:\PrecSpace{\sigma}\to\SubTarget{\sigma}.
    \]
    Define the precursor restriction functional
    \begin{equation}
    \PrecFunctional{\sigma}(p)
    \equiv
    \SubPair{\sigma}\bigl(\PrecSection{\sigma}(p),\PrecSection{\sigma}(p)\bigr),
    \label{eq:precfunctional}
    \end{equation}
    together with the exact precursor shell
    \begin{equation}
    \PrecShellSet{\sigma}
    \equiv
    \bigl(\PrecSection{\sigma}\bigr)^{-1}(0)\cap\PrecBody{\sigma}.
    \label{eq:precshell}
    \end{equation}
    Define the induced lifted restriction section
    \begin{equation}
    \LiftSection{\sigma}
    \equiv
    \PrecSection{\sigma}\circ\PrecLiftMin{\sigma},
    \label{eq:inducedliftsection}
    \end{equation}
    and the induced lifted restriction functional
    \begin{equation}
    \LiftFunctional{\sigma}(\ell)
    \equiv
    \SubPair{\sigma}\bigl(\LiftSection{\sigma}(\ell),\LiftSection{\sigma}(\ell)\bigr)
    =
    \PrecFunctional{\sigma}\bigl(\PrecLiftMin{\sigma}(\ell)\bigr).
    \label{eq:inducedliftfunctional}
    \end{equation}
    Assume that, for each \(\ell\in\LiftBody{\sigma}\), \(\PrecLiftMin{\sigma}(\ell)\) is the unique minimizer of \(\PrecFunctional{\sigma}\) on the fiber
    \[
    \bigl(\PrecProj{\sigma}\bigr)^{-1}(\ell)\cap\PrecBody{\sigma},
    \]
    and, for each \(m\in\SubBody{\sigma}\), \(\PrecSubMin{\sigma}(m)\) is the unique minimizer of \(\PrecFunctional{\sigma}\) on the composite fiber
    \[
    \bigl(\LiftProj{\sigma}\circ\PrecProj{\sigma}\bigr)^{-1}(m)\cap\PrecBody{\sigma}.
    \]
    Assume that \(\LiftSelect{\sigma}(m)\) is therefore the unique minimizer of \(\LiftFunctional{\sigma}\) on
    \[
    \bigl(\LiftProj{\sigma}\bigr)^{-1}(m)\cap\LiftBody{\sigma},
    \]
    and that there exists \(\mu_{\sigma}^{\mathrm{fib}}>0\) such that
    \begin{equation}
    \LiftFunctional{\sigma}(\ell)
    \ge
    \LiftFunctional{\sigma}\bigl(\LiftSelect{\sigma}(m)\bigr)
    +
    \mu_{\sigma}^{\mathrm{fib}}\,
    \operatorname{dist}\bigl(\ell,\LiftSelect{\sigma}(m)\bigr)^{2}
    \label{eq:precfibergap}
    \end{equation}
    for every \(m\in\SubBody{\sigma}\) and every \(\ell\in\bigl(\LiftProj{\sigma}\bigr)^{-1}(m)\cap\LiftBody{\sigma}\). Assume there exists \(\mu_{\sigma}^{\mathrm{sub}}>0\) such that
    \begin{equation}
    \LiftFunctional{\sigma}\bigl(\LiftSelect{\sigma}(m)\bigr)
    \ge
    \mu_{\sigma}^{\mathrm{sub}}\,
    \operatorname{dist}\bigl(m,\LiftProj{\sigma}(\PrecProj{\sigma}(\PrecShellSet{\sigma}))\bigr)^{2}
    \qquad
    \text{for all }m\in\SubBody{\sigma}.
    \label{eq:precprojectedgap}
    \end{equation}
    \item Affine lifted observer-data and shell-response readouts
    \[
    \LiftData{\sigma}:\LiftSpace{\sigma}\to\DataSpace{\sigma},
    \qquad
    \LiftShellMap{\sigma}:\LiftSpace{\sigma}\to\AmbientTarget{\sigma},
    \]
    and affine precursor observer-data and shell-response readouts
    \[
    \PrecData{\sigma}:\PrecSpace{\sigma}\to\DataSpace{\sigma},
    \qquad
    \PrecShellMap{\sigma}:\PrecSpace{\sigma}\to\AmbientTarget{\sigma},
    \]
    satisfying
    \begin{equation}
    \PrecData{\sigma}
    =
    \LiftData{\sigma}\circ\PrecProj{\sigma},
    \qquad
    \PrecShellMap{\sigma}
    =
    \LiftShellMap{\sigma}\circ\PrecProj{\sigma}.
    \label{eq:precfactor}
    \end{equation}
    There is also a finite-dimensional Euclidean support target
    \[
    \ResTarget{\sigma},
    \]
    a positive-definite symmetric support pairing
    \[
    \ResPair{\sigma}:
    \ResTarget{\sigma}\times\ResTarget{\sigma}\to\mathbb R,
    \]
    a finite-dimensional real hidden target
    \[
    \EqnTarget{\sigma},
    \]
    and a positive-definite symmetric hidden pairing
    \[
    \EqnPair{\sigma}:
    \EqnTarget{\sigma}\times\EqnTarget{\sigma}\to\mathbb R.
    \]
    There is a compact symmetry group
    \[
    \PrecSym{\sigma}
    \]
    and an involution
    \[
    \PrecTime{\sigma},
    \]
    acting orthogonally on \(\PrecSpace{\sigma}\), \(\LiftSpace{\sigma}\), \(\SubSpace{\sigma}\), \(\DataSpace{\sigma}\), \(\AmbientTarget{\sigma}\), and \(\ResTarget{\sigma}\), and by \(\EqnPair{\sigma}\)-isometries on \(\EqnTarget{\sigma}\), preserving \(\PrecBody{\sigma}\), \(\LiftBody{\sigma}\), \(\PrecProj{\sigma}\), \(\LiftProj{\sigma}\), \(\PrecLiftMin{\sigma}\), \(\PrecSubMin{\sigma}\), \(\PrecSection{\sigma}\), \(\LiftData{\sigma}\), \(\LiftShellMap{\sigma}\), \(\PrecData{\sigma}\), and \(\PrecShellMap{\sigma}\).
    \item The restriction of \(\PrecProj{\sigma}\) to \(\PrecShellSet{\sigma}\) is a diffeomorphism onto its image
    \[
    \LiftShellSet{\sigma}
    \equiv
    \PrecProj{\sigma}\bigl(\PrecShellSet{\sigma}\bigr)
    \subset\LiftBody{\sigma}.
    \]
    There is a Euclidean precursor shell chart
    \[
    \PrecChart{\sigma}:\PrecShellSet{\sigma}\to\ResSpace{\sigma}
    \]
    whose shell-body image
    \[
    \ResBody{\sigma}
    \equiv
    \PrecChart{\sigma}\bigl(\PrecShellSet{\sigma}\bigr)
    \]
    is compact convex with nonempty interior. There exist affine maps
    \[
    \ResData{\sigma}:\ResSpace{\sigma}\to\DataSpace{\sigma},
    \qquad
    \ResShellMap{\sigma}:\ResSpace{\sigma}\to\AmbientTarget{\sigma},
    \]
    and a smooth support section
    \[
    \ResSection{\sigma}:\ResSpace{\sigma}\to\ResTarget{\sigma}
    \]
    such that
    \begin{equation}
    \ResData{\sigma}\circ\PrecChart{\sigma}
    =
    \PrecData{\sigma}\big|_{\PrecShellSet{\sigma}},
    \qquad
    \ResShellMap{\sigma}\circ\PrecChart{\sigma}
    =
    \PrecShellMap{\sigma}\big|_{\PrecShellSet{\sigma}}.
    \label{eq:precchartreadouts}
    \end{equation}
    There is a smooth precursor shell-internal support recorder
    \[
    \PrecSupport{\sigma}:\PrecShellSet{\sigma}\to\ResTarget{\sigma}
    \]
    satisfying
    \begin{equation}
    \ResSection{\sigma}\circ\PrecChart{\sigma}
    =
    \PrecSupport{\sigma}.
    \label{eq:precsupporttransport}
    \end{equation}
    Define the exact precursor support shell
    \begin{equation}
    \PrecSupportShell{\sigma}
    \equiv
    \bigl(\PrecSupport{\sigma}\bigr)^{-1}(0)\cap\PrecShellSet{\sigma},
    \label{eq:precsupportshell}
    \end{equation}
    and assume
    \begin{equation}
    \ResShellSet{\sigma}
    \equiv
    \bigl(\ResSection{\sigma}\bigr)^{-1}(0)\cap\ResBody{\sigma}
    =
    \PrecChart{\sigma}\bigl(\PrecSupportShell{\sigma}\bigr).
    \label{eq:prectransportedsupportshell}
    \end{equation}
    Assume there exists \(\nu_{\sigma}^{\mathrm{sup}}>0\) such that
    \begin{equation}
    \ResPair{\sigma}\bigl(\ResSection{\sigma}(u),\ResSection{\sigma}(u)\bigr)
    \ge
    \nu_{\sigma}^{\mathrm{sup}}\,
    \operatorname{dist}\bigl(u,\ResShellSet{\sigma}\bigr)^{2}
    \qquad
    \text{for all }u\in\ResBody{\sigma}.
    \label{eq:precnormal}
    \end{equation}
    Assume the restriction of \(\PrecProj{\sigma}\) to \(\PrecSupportShell{\sigma}\) is a diffeomorphism onto its image
    \[
    \LiftSupportShell{\sigma}
    \equiv
    \PrecProj{\sigma}\bigl(\PrecSupportShell{\sigma}\bigr).
    \]
    The transported symmetry and time-reversal actions act linearly on \(\ResSpace{\sigma}\), preserving \(\ResBody{\sigma}\), \(\ResShellSet{\sigma}\), \(\ResSection{\sigma}\), \(\ResData{\sigma}\), and \(\ResShellMap{\sigma}\).
    \item A closed embedded precursor support zero core
    \[
    \PrecZeroCore{\sigma}\subset\PrecSupportShell{\sigma},
    \]
    a finite-dimensional Euclidean hidden-seed space
    \[
    \SeedHidden{\sigma},
    \]
    and a compact convex hidden-seed body
    \[
    \SeedHiddenBody{\sigma}\subset\SeedHidden{\sigma}
    \]
    with
    \[
    0\in\operatorname{int}\SeedHiddenBody{\sigma}.
    \]
    There is a Euclidean precursor support-shell chart
    \[
    \PrecEqChart{\sigma}:\PrecSupportShell{\sigma}\to\EqnSpace{\sigma},
    \]
    an affine chart on the precursor support zero core
    \[
    \PrecZeroChart{\sigma}:\PrecZeroCore{\sigma}\to\EqnZero{\sigma},
    \]
    and a linear isometric embedding
    \[
    \HiddenChart{\sigma}:\SeedHidden{\sigma}\to\EqnHidden{\sigma}
    \]
    such that
    \[
    \EqnSpace{\sigma}
    =
    \EqnZero{\sigma}\oplus\EqnHidden{\sigma}
    \]
    is an orthogonal direct sum, the set
    \[
    \ZeroBody{\sigma}
    \equiv
    \PrecZeroChart{\sigma}\bigl(\PrecZeroCore{\sigma}\bigr)
    \]
    is compact convex with nonempty interior in \(\EqnZero{\sigma}\), the shell image
    \[
    \EqnBody{\sigma}
    \equiv
    \PrecEqChart{\sigma}\bigl(\PrecSupportShell{\sigma}\bigr)
    \]
    has nonempty interior in \(\EqnSpace{\sigma}\), and there is a diffeomorphism
    \begin{equation}
    \PrecSplit{\sigma}:
    \PrecZeroCore{\sigma}\times\SeedHiddenBody{\sigma}
    \xrightarrow{\cong}
    \PrecSupportShell{\sigma}
    \label{eq:precsplit}
    \end{equation}
    satisfying
    \begin{equation}
    \PrecEqChart{\sigma}\bigl(\PrecSplit{\sigma}(z,\eta)\bigr)
    =
    \PrecZeroChart{\sigma}(z)+\HiddenChart{\sigma}(\eta).
    \label{eq:precchartsplit}
    \end{equation}
    The precursor shell readouts become affine in the precursor support-shell chart: there exist affine maps
    \[
    \EqnData{\sigma}:\EqnSpace{\sigma}\to\DataSpace{\sigma},
    \qquad
    \EqnShell{\sigma}:\EqnSpace{\sigma}\to\AmbientTarget{\sigma},
    \]
    such that
    \begin{equation}
    \EqnData{\sigma}\circ\PrecEqChart{\sigma}
    =
    \PrecData{\sigma}\big|_{\PrecSupportShell{\sigma}},
    \qquad
    \EqnShell{\sigma}\circ\PrecEqChart{\sigma}
    =
    \PrecShellMap{\sigma}\big|_{\PrecSupportShell{\sigma}}.
    \label{eq:preceqnchartreadouts}
    \end{equation}
    Assume the restriction of \(\PrecProj{\sigma}\) to \(\PrecZeroCore{\sigma}\) is a diffeomorphism onto its image
    \[
    \LiftZeroCore{\sigma}
    \equiv
    \PrecProj{\sigma}\bigl(\PrecZeroCore{\sigma}\bigr).
    \]
    The transported symmetry and time-reversal actions act linearly on \(\EqnSpace{\sigma}\), preserving \(\EqnZero{\sigma}\), \(\EqnHidden{\sigma}\), and \(\EqnBody{\sigma}\).
    \item A linear hidden charge map
    \[
    \PrecCharge{\sigma}:\SeedHidden{\sigma}\to\EqnTarget{\sigma}
    \]
    such that there exists \(\lambda_{\sigma}^{\mathrm{eqn}}>0\) with
    \begin{equation}
    \EqnPair{\sigma}\bigl(\PrecCharge{\sigma}\eta,\PrecCharge{\sigma}\eta\bigr)
    \ge
    \lambda_{\sigma}^{\mathrm{eqn}}\,
    \|\eta\|^{2}
    \qquad
    \text{for all }\eta\in\SeedHidden{\sigma}.
    \label{eq:precchargegap}
    \end{equation}
    Define the precursor support-shell hidden recorder by
    \begin{equation}
    \PrecResidual{\sigma}\bigl(\PrecSplit{\sigma}(z,\eta)\bigr)
    \equiv
    \PrecCharge{\sigma}(\eta).
    \label{eq:precresidual}
    \end{equation}
    \item A smooth embedding
    \[
    \jmath_{\sigma}:X_{\sigma}\hookrightarrow\VisibleAffine{\sigma}
    \]
    and a smooth observer-compatible exact precursor zero-response realization
    \[
    \PrecEmbed{\sigma}:\ForSpace{\sigma}\hookrightarrow\PrecZeroCore{\sigma}
    \]
    whose image is exactly
    \[
    \PrecZeroCore{\sigma}
    \cap
    \bigl(\PrecShellMap{\sigma}\bigr)^{-1}(0),
    \]
    and such that
    \begin{equation}
    \PrecData{\sigma}\bigl(\PrecEmbed{\sigma}(w)\bigr)
    =
    \bigl(
    \jmath_{\sigma}(\Reduction{\sigma}(w)),
    \CurrentCarrier{\sigma}(w),
    \EqCarrier{\sigma}(w)
    \bigr)
    \label{eq:preccompat}
    \end{equation}
    for every \(w\in\ForSpace{\sigma}\).
    \item Fixed-data solvability on the precursor support zero core:
    \begin{equation}
    \PrecData{\sigma}\bigl(\PrecZeroCore{\sigma}\bigr)
    =
    \PrecData{\sigma}\Bigl(
    \PrecZeroCore{\sigma}\cap
    \bigl(\PrecShellMap{\sigma}\bigr)^{-1}(0)
    \Bigr).
    \label{eq:precsolvability}
    \end{equation}
    \item The inert precursor support-zero-core kernel
    \[
    \PrecKernel{\sigma}
    \equiv
    \ker\bigl(D\PrecData{\sigma}\big|_{\PrecZeroCore{\sigma}}\bigr)
    \cap
    \ker\bigl(D\PrecShellMap{\sigma}\big|_{\PrecZeroCore{\sigma}}\bigr)
    \]
    and the orthogonal projector
    \[
    \PrecStateComp{\sigma}:
    \EqnZero{\sigma}\to
    \bigl(\PrecKernelCh{\sigma}\bigr)^{\perp}\cap\EqnZero{\sigma},
    \qquad
    \PrecKernelCh{\sigma}
    \equiv
    D\PrecZeroChart{\sigma}\bigl(\PrecKernel{\sigma}\bigr).
    \]
    There exists \(\kappa_{\sigma}^{\mathrm{eqn}}>0\) such that for every
    \[
    w\in\ForSpace{\sigma},
    \qquad
    \delta p\in T_{\PrecEmbed{\sigma}(w)}\PrecSupportShell{\sigma}
    \]
    satisfying
    \[
    D\PrecData{\sigma}\,\delta p=0,
    \]
    if
    \[
    D\PrecEqChart{\sigma}(\delta p)=\delta z+\delta\eta,
    \qquad
    \delta z\in\EqnZero{\sigma},
    \quad
    \delta\eta\in\EqnHidden{\sigma},
    \]
    then
    \begin{equation}
    \begin{aligned}
    \|D\PrecShellMap{\sigma}\,\delta p\|^{2}
    &+
    \EqnPair{\sigma}\bigl(
    \PrecCharge{\sigma}\bigl((\HiddenChart{\sigma})^{-1}\delta\eta\bigr),\\
    &\qquad
    \PrecCharge{\sigma}\bigl((\HiddenChart{\sigma})^{-1}\delta\eta\bigr)
    \bigr)
    \\
    &\ge
    \kappa_{\sigma}^{\mathrm{eqn}}
    \left(
    \|\PrecStateComp{\sigma}(\delta z)\|^{2}
    +
    \|\delta\eta\|^{2}
    \right).
    \end{aligned}
    \label{eq:preccoercive}
    \end{equation}
\end{enumerate}
\end{definition}

\begin{remark}
Definition \ref{def:prec} is deliberately deeper than the shell-generating lift package but narrower than a claim of completed substrate microphysics. The lift layer is not discarded. Instead it is shown to arise as the projected exact-shell image of a still more primitive precursor layer whose exact shell, support shell, zero-core / hidden-seed split, compatibility realization, solvability statement, and coercive control all first live at the precursor level.
\end{remark}

\section{Canonical Shell-Generating Lift Package}

\begin{proposition}[The shell-generating substrate lift package is produced canonically]
\label{prop:lift}
Assume Definition \ref{def:prec}. Define, for each sector \(\sigma\in\{\Car,\Cur,\Hom\}\),
\begin{equation}
\LiftSection{\sigma}
\equiv
\PrecSection{\sigma}\circ\PrecLiftMin{\sigma},
\label{eq:propinducedliftsection}
\end{equation}
the lifted shell chart
\begin{equation}
\LiftChart{\sigma}
\equiv
\PrecChart{\sigma}\circ
\bigl(\PrecProj{\sigma}\big|_{\PrecShellSet{\sigma}}\bigr)^{-1},
\label{eq:propinducedliftchart}
\end{equation}
the lifted shell-internal support recorder
\begin{equation}
\LiftSupport{\sigma}
\equiv
\PrecSupport{\sigma}\circ
\bigl(\PrecProj{\sigma}\big|_{\PrecShellSet{\sigma}}\bigr)^{-1},
\label{eq:propinducedliftsupport}
\end{equation}
the projected support shell and projected support zero core
\begin{equation}
\LiftSupportShell{\sigma}
\equiv
\PrecProj{\sigma}\bigl(\PrecSupportShell{\sigma}\bigr),
\qquad
\LiftZeroCore{\sigma}
\equiv
\PrecProj{\sigma}\bigl(\PrecZeroCore{\sigma}\bigr),
\label{eq:propinducedshells}
\end{equation}
the support-shell and zero-core charts
\begin{equation}
\LiftEqChart{\sigma}
\equiv
\PrecEqChart{\sigma}\circ
\bigl(\PrecProj{\sigma}\big|_{\PrecSupportShell{\sigma}}\bigr)^{-1},
\qquad
\LiftZeroChart{\sigma}
\equiv
\PrecZeroChart{\sigma}\circ
\bigl(\PrecProj{\sigma}\big|_{\PrecZeroCore{\sigma}}\bigr)^{-1},
\label{eq:propinducedcharts}
\end{equation}
the projected split map
\begin{equation}
\LiftSplit{\sigma}(z,\eta)
\equiv
\PrecProj{\sigma}\Bigl(
\PrecSplit{\sigma}\bigl(
(\PrecProj{\sigma}\big|_{\PrecZeroCore{\sigma}})^{-1}(z),
\eta
\bigr)
\Bigr),
\label{eq:propinducedsplit}
\end{equation}
the hidden charge and hidden recorder
\begin{equation}
\LiftCharge{\sigma}
\equiv
\PrecCharge{\sigma},
\qquad
\LiftResidual{\sigma}\bigl(\LiftSplit{\sigma}(z,\eta)\bigr)
\equiv
\PrecCharge{\sigma}(\eta),
\label{eq:propinducedresidual}
\end{equation}
and the exact zero-response realization
\begin{equation}
\LiftEmbed{\sigma}
\equiv
\PrecProj{\sigma}\circ\PrecEmbed{\sigma}.
\label{eq:propinducedembed}
\end{equation}
Then the resulting data satisfy exactly the shell-generating substrate lift dynamics package of Definition \ref{def:lift}.
\end{proposition}

\begin{proof}
We verify the items of Definition \ref{def:lift} in order.

For item (1), the lifted state space \(\LiftSpace{\sigma}\), the compact convex lifted body \(\LiftBody{\sigma}\), the affine surjection \(\LiftProj{\sigma}\), and the projected substrate body \(\SubBody{\sigma}\) are already part of Definition \ref{def:prec}(1). Equation \eqref{eq:inducedliftselect} gives the smooth selection \(\LiftSelect{\sigma}\), and
\[
\LiftProj{\sigma}\circ\LiftSelect{\sigma}
=
\LiftProj{\sigma}\circ\PrecProj{\sigma}\circ\PrecSubMin{\sigma}
=
\mathrm{id}_{\SubBody{\sigma}}
\]
by Definition \ref{def:prec}(1).

For item (2), equation \eqref{eq:propinducedliftsection} gives the lifted restriction section, and \eqref{eq:inducedliftfunctional} gives its associated lifted restriction functional. If \(\ell\in\PrecProj{\sigma}(\PrecShellSet{\sigma})\), choose \(p\in\PrecShellSet{\sigma}\) with \(\PrecProj{\sigma}(p)=\ell\). Then \(\PrecFunctional{\sigma}(p)=0\), so by the precursor lift-fiber minimality of Definition \ref{def:prec}(2),
\[
\LiftFunctional{\sigma}(\ell)
=
\PrecFunctional{\sigma}\bigl(\PrecLiftMin{\sigma}(\ell)\bigr)
=0.
\]
Conversely, if \(\LiftFunctional{\sigma}(\ell)=0\), then
\[
\PrecFunctional{\sigma}\bigl(\PrecLiftMin{\sigma}(\ell)\bigr)=0,
\]
so \(\PrecLiftMin{\sigma}(\ell)\in\PrecShellSet{\sigma}\) and therefore \(\ell\in\PrecProj{\sigma}(\PrecShellSet{\sigma})\). Hence
\[
\LiftShellSet{\sigma}
=
\bigl(\LiftSection{\sigma}\bigr)^{-1}(0)\cap\LiftBody{\sigma}
=
\PrecProj{\sigma}\bigl(\PrecShellSet{\sigma}\bigr).
\]
The fiber-minimizing property of \(\LiftSelect{\sigma}\) and the estimate \eqref{eq:fibergap} are exactly the assumptions \eqref{eq:precfibergap} of Definition \ref{def:prec}(2), while \eqref{eq:projectedgap} is exactly \eqref{eq:precprojectedgap}.

For item (3), the affine readouts \(\LiftData{\sigma}\) and \(\LiftShellMap{\sigma}\), the support target, support pairing, hidden target, hidden pairing, and the symmetry/time-reversal actions on \(\LiftSpace{\sigma}\) are already part of Definition \ref{def:prec}(3). Those same actions give the induced \(\LiftSym{\sigma}\) and \(\LiftTime{\sigma}\). Because Definition \ref{def:prec}(3) says they preserve \(\PrecProj{\sigma}\), \(\LiftProj{\sigma}\), \(\PrecSubMin{\sigma}\), \(\PrecLiftMin{\sigma}\), \(\PrecSection{\sigma}\), \(\LiftData{\sigma}\), and \(\LiftShellMap{\sigma}\), they preserve \(\LiftBody{\sigma}\), \(\LiftSelect{\sigma}\), \(\LiftSection{\sigma}\), and the full lifted readout package required in Definition \ref{def:lift}(3).

For item (4), the restriction \(\PrecProj{\sigma}\big|_{\PrecShellSet{\sigma}}\) is a diffeomorphism onto \(\LiftShellSet{\sigma}\) by Definition \ref{def:prec}(4), so \eqref{eq:propinducedliftchart} defines the lifted shell chart. Its image is exactly
\[
\ResBody{\sigma}
=
\PrecChart{\sigma}\bigl(\PrecShellSet{\sigma}\bigr).
\]
Using \eqref{eq:precfactor} and \eqref{eq:precchartreadouts},
\[
\ResData{\sigma}\circ\LiftChart{\sigma}
=
\PrecData{\sigma}\circ
\bigl(\PrecProj{\sigma}\big|_{\PrecShellSet{\sigma}}\bigr)^{-1}
=
\LiftData{\sigma}\big|_{\LiftShellSet{\sigma}},
\]
and similarly
\[
\ResShellMap{\sigma}\circ\LiftChart{\sigma}
=
\LiftShellMap{\sigma}\big|_{\LiftShellSet{\sigma}}.
\]
Equation \eqref{eq:propinducedliftsupport} defines the lifted shell-internal support recorder, and \eqref{eq:precsupporttransport} becomes exactly \eqref{eq:liftsupporttransport}. Because the restriction of \(\PrecProj{\sigma}\) to \(\PrecSupportShell{\sigma}\) is a diffeomorphism onto \(\LiftSupportShell{\sigma}\), one has
\[
\ResShellSet{\sigma}
=
\PrecChart{\sigma}\bigl(\PrecSupportShell{\sigma}\bigr)
=
\LiftChart{\sigma}\bigl(\LiftSupportShell{\sigma}\bigr),
\]
which is \eqref{eq:lifttransportedsupportshell}. The normal shell stability \eqref{eq:liftnormal} is exactly \eqref{eq:precnormal}. The transported symmetry and time-reversal actions on \(\ResSpace{\sigma}\) are the same ones already supplied in Definition \ref{def:prec}(4).

For item (5), the projected support zero core is \(\LiftZeroCore{\sigma}\), and the restrictions of \(\PrecProj{\sigma}\big|_{\PrecSupportShell{\sigma}}\) and \(\PrecProj{\sigma}\big|_{\PrecZeroCore{\sigma}}\) are diffeomorphisms onto \(\LiftSupportShell{\sigma}\) and \(\LiftZeroCore{\sigma}\), respectively. Therefore \eqref{eq:propinducedcharts} defines the required charts. The orthogonal splitting
\[
\EqnSpace{\sigma}=\EqnZero{\sigma}\oplus\EqnHidden{\sigma},
\]
the compact convex bodies \(\ZeroBody{\sigma}\) and \(\EqnBody{\sigma}\), and the isometric hidden chart \(\HiddenChart{\sigma}\) are already part of Definition \ref{def:prec}(5). The projected split map \eqref{eq:propinducedsplit} is a diffeomorphism because \(\PrecSplit{\sigma}\) is one and \(\PrecProj{\sigma}\big|_{\PrecSupportShell{\sigma}}\) is one. Moreover, if
\[
p_{z}
\equiv
\bigl(\PrecProj{\sigma}\big|_{\PrecZeroCore{\sigma}}\bigr)^{-1}(z)
\in
\PrecZeroCore{\sigma},
\]
then
\[
\LiftEqChart{\sigma}\bigl(\LiftSplit{\sigma}(z,\eta)\bigr)
=
\PrecEqChart{\sigma}\bigl(\PrecSplit{\sigma}(p_{z},\eta)\bigr)
=
\PrecZeroChart{\sigma}(p_{z})+\HiddenChart{\sigma}(\eta)
=
\LiftZeroChart{\sigma}(z)+\HiddenChart{\sigma}(\eta),
\]
which is exactly \eqref{eq:liftchartsplit}. The equation-space readouts \(\EqnData{\sigma}\) and \(\EqnShell{\sigma}\) are precisely the affine maps already supplied in Definition \ref{def:prec}(5). Using \eqref{eq:precfactor}, equation \eqref{eq:preceqnchartreadouts} becomes \eqref{eq:lifteqnchartreadouts}.

For item (6), \eqref{eq:propinducedresidual} defines the hidden charge and hidden recorder, and the positive hidden gap \eqref{eq:liftchargegap} is exactly \eqref{eq:precchargegap}.

For item (7), \eqref{eq:propinducedembed} defines the exact lifted zero-response realization. Since \(\PrecEmbed{\sigma}\) lands in \(\PrecZeroCore{\sigma}\cap(\PrecShellMap{\sigma})^{-1}(0)\), its projected image lands in \(\LiftZeroCore{\sigma}\cap(\LiftShellMap{\sigma})^{-1}(0)\) by \eqref{eq:precfactor}. Because \(\PrecProj{\sigma}\big|_{\PrecZeroCore{\sigma}}\) is a diffeomorphism onto \(\LiftZeroCore{\sigma}\), the image is exactly that zero-response locus. Equation \eqref{eq:liftcompat} follows immediately from \eqref{eq:preccompat}, \eqref{eq:precfactor}, and \eqref{eq:propinducedembed}.

For item (8), fixed-data solvability \eqref{eq:liftsolvability} is the projection of \eqref{eq:precsolvability} through the zero-core diffeomorphism \(\PrecProj{\sigma}\big|_{\PrecZeroCore{\sigma}}\) together with the factorization \eqref{eq:precfactor}.

For item (9), define
\[
\LiftKernel{\sigma}
\equiv
\ker\bigl(D\LiftData{\sigma}\big|_{\LiftZeroCore{\sigma}}\bigr)
\cap
\ker\bigl(D\LiftShellMap{\sigma}\big|_{\LiftZeroCore{\sigma}}\bigr),
\qquad
\LiftKernelCh{\sigma}
\equiv
D\LiftZeroChart{\sigma}\bigl(\LiftKernel{\sigma}\bigr),
\]
and let \(\LiftStateComp{\sigma}\) be the orthogonal projector onto \((\LiftKernelCh{\sigma})^{\perp}\cap\EqnZero{\sigma}\). Take \(w\in\ForSpace{\sigma}\) and \(\delta \ell\in T_{\LiftEmbed{\sigma}(w)}\LiftSupportShell{\sigma}\) satisfying \(D\LiftData{\sigma}\,\delta \ell=0\). Set
\[
\delta p
\equiv
D\bigl(\PrecProj{\sigma}\big|_{\PrecSupportShell{\sigma}}\bigr)^{-1}(\delta \ell)
\in
T_{\PrecEmbed{\sigma}(w)}\PrecSupportShell{\sigma}.
\]
Then \(D\PrecData{\sigma}\,\delta p=0\) by \eqref{eq:precfactor}. Writing
\[
D\LiftEqChart{\sigma}(\delta \ell)=\delta z+\delta\eta,
\qquad
\delta z\in\EqnZero{\sigma},
\quad
\delta\eta\in\EqnHidden{\sigma},
\]
equation \eqref{eq:preccoercive} becomes exactly \eqref{eq:liftcoercive}. This completes the verification of Definition \ref{def:lift}.
\end{proof}

\begin{remark}
Proposition \ref{prop:lift} gives more than another same-output relay. Every constituent of the shell-generating substrate lift package is now an explicit projected or transported image of deeper precursor data. The lift layer has therefore lost its status as a free insertion on the active line.
\end{remark}

\section{Necessity and Uniqueness at Lift Scope}

\begin{proposition}[No extra benchmark-facing freedom remains at lift scope]
\label{prop:unique}
Assume the lift-generating substrate precursor dynamics package of Definition \ref{def:prec}. Then for each sector \(\sigma\):
\begin{enumerate}
    \item the lifted state space and lifted body, the substrate projection and fiber-minimizing lifted selection, the lifted restriction section and exact lifted shell, the affine lifted observer-data and shell-response readouts, the lifted shell chart and shell-internal support recorder, the lifted support shell and lifted support zero core, the charted zero-core / hidden-seed split, the hidden charge and hidden recorder, the symmetry and time-reversal structure, the exact lifted zero-response realization, the fixed-data solvability statement, the inert kernel/projector, and the lifted coercive package are all canonically induced by the precursor input;
    \item any other shell-generating lift presentation compatible with that same precursor input differs from the canonical one only by the obvious precursor-faithful identification preserving the same lifted fibers, the same charted support shell and support zero core, the same zero-hidden splitting and hidden quadratic form, and the same benchmark shell data.
\end{enumerate}
\end{proposition}

\begin{proof}
Item (1) is Proposition \ref{prop:lift}. For item (2), the precursor input already fixes the lifted body as \(\PrecProj{\sigma}(\PrecBody{\sigma})\), the substrate projection \(\LiftProj{\sigma}\), the fiber-minimizing lift \(\LiftSelect{\sigma}=\PrecProj{\sigma}\circ\PrecSubMin{\sigma}\), the lifted restriction section \(\LiftSection{\sigma}=\PrecSection{\sigma}\circ\PrecLiftMin{\sigma}\), the exact lifted shell \(\LiftShellSet{\sigma}=\PrecProj{\sigma}(\PrecShellSet{\sigma})\), the lifted shell chart and support recorder through \eqref{eq:propinducedliftchart} and \eqref{eq:propinducedliftsupport}, the lifted support shell and support zero core through \eqref{eq:propinducedshells}, the lifted support-shell and zero-core charts through \eqref{eq:propinducedcharts}, the charted split through \eqref{eq:propinducedsplit}, the hidden quadratic form through \(\LiftCharge{\sigma}=\PrecCharge{\sigma}\) and \(\EqnPair{\sigma}\), the exact zero-response realization through \eqref{eq:propinducedembed}, and the solvability and coercive statements through the transported precursor identities. Any compatible lift presentation must therefore preserve precisely those same transported objects. That leaves only the obvious precursor-faithful identification respecting the already fixed lifted fibers, the same charted support shell, the same zero/hidden splitting and hidden quadratic form, and the same benchmark shell data.
\end{proof}

\begin{corollary}[The lift package is forced by precursor dynamics]
\label{cor:forced}
Within the class of lift-generating substrate precursor dynamics packages, the shell-generating substrate lift dynamics are necessary and unique at benchmark-facing scope.
\end{corollary}

\begin{proof}
This is the combined content of Proposition \ref{prop:lift} and Proposition \ref{prop:unique}.
\end{proof}

\section{Benchmark Preservation}

\begin{theorem}[The one-metric benchmark package is unchanged]
\label{thm:outputs}
Assume Definition \ref{def:prec}. Then the benchmark output statements of Definition \ref{def:fixed} remain unchanged:
\[
\MetricOut(\mathbf w)=\CompletedMetric,
\qquad
\MatterOut(\mathbf w,\psi)
=
T_{\mu\nu}^{\mathrm{matter}}[\CompletedMetric,\psi]
\]
for every \(\mathbf w\in\PrimShell\) and every matter configuration \(\psi\in\Psi\). Consequently the lift-forcing result introduces neither a second operative metric nor an enlarged low-energy matter law.
\end{theorem}

\begin{proof}
The present note changes only the upstream benchmark-facing status of the shell-generating substrate lift package. It does not alter the primitive shell \(\PrimShell\), the recorded metric map \(\MetricOut\), or the recorded matter map \(\MatterOut\). Hence \eqref{eq:fixedoutputs} continues to hold exactly.
\end{proof}

\section{Main Theorem}

\begin{theorem}[The remaining burden moves below lift scope]
\label{thm:main}
Assume the fixed primitive continuation data of Definition \ref{def:fixed} and, for each sector \(\sigma\in\{\Car,\Cur,\Hom\}\), lift-generating substrate precursor dynamics in the sense of Definition \ref{def:prec}. Then:
\begin{enumerate}
    \item the shell-generating substrate lift dynamics are canonically induced by that precursor package;
    \item the lift package is necessary and unique there up to the obvious precursor-faithful identification preserving the same lifted fibers, the same charted support shell and support zero core, the same zero-hidden splitting and hidden quadratic form, and the same benchmark shell data;
    \item the one operative metric remains exactly \(\CompletedMetric\), the ordinary matter route remains exactly the standard one through that metric, there is no second operative metric, and the low-energy matter law is not enlarged;
    \item consequently the benchmark-facing burden after the shell-restriction forcing theorem sharpens once more: the shell-generating substrate lift package is no longer left as an independently inserted stopping point, and the honest burden now lies below lift scope at the precursor layer itself.
\end{enumerate}
\end{theorem}

\begin{proof}
Item (1) is Proposition \ref{prop:lift}. Item (2) is Corollary \ref{cor:forced}. Item (3) is Theorem \ref{thm:outputs}. Item (4) is the routing consequence of the previous items.
\end{proof}

\begin{corollary}[What remains open after this theorem]
\label{cor:next}
After Theorem \ref{thm:main}, the remaining honest burden is no longer to justify the shell-generating substrate lift dynamics themselves on the active primitive-reservoir observer-reduction line. The next honest burden is sharper again: first try to derive or uniquely force the lift-generating substrate precursor dynamics from still more primitive explicit substrate structure in a way that actually changes benchmark-facing status, or else prove a genuinely smaller theorem-level collapse, sharper necessity result, or sharper uniqueness result strictly inside the shell-generating lift package while preserving the same benchmark one-metric package and claim boundary.
\end{corollary}

\begin{proof}
Theorem \ref{thm:main} converts the lift package from the remaining benchmark-facing stopping point into a canonical output of the precursor class. What remains open is therefore the precursor layer itself, not the lift package it now forces.
\end{proof}

\section{Conclusion}

The positive result of this note is deliberately narrower than a completed final microphysical derivation and stronger than another lift-level relay. The shell-restriction forcing theorem had already shown that the split zero-hidden selector-shell core carries no extra benchmark-facing freedom once the shell-restriction layer is fixed. The present note now makes the next admissible upstream forcing move below that frontier: it shows that the shell-generating substrate lift dynamics themselves are a canonical and unique output of lift-generating substrate precursor dynamics.

That gain is exactly what the routing layer required. The lift package is no longer left as an independently inserted stopping point. Its lifted body and substrate projection, fiber-minimizing lift, lifted restriction section, exact lifted shell, lifted shell/support charts, lifted support zero core, zero-hidden split, hidden quadratic form, exact zero-response realization, solvability statement, inert-kernel projector, and coercive control are all already fixed once the precursor input is fixed. At current scope, any compatible lift presentation differs only by the obvious precursor-faithful identification preserving those same benchmark-facing data.

The benchmark boundary remains unchanged. The same primitive shell remains the shell carrying the observer outputs, so the one operative metric is still exactly \(\CompletedMetric\), the ordinary matter route is unchanged, there is no second operative metric, and the low-energy matter law is not enlarged. Nothing here upgrades adoption to derivation or licenses stronger first-principles promotion language.

The open burden therefore moves below lift scope. The next benchmark-facing move should not spend itself on another shell-restriction relay, another split-core relay, another selector-law restatement, another combined-response restatement, another ambient-package restatement, another trace-core restatement, another completion-core restatement, another exact-shell affine nucleus restatement, or another same lift relay. The honest next step is to go one level lower and address the precursor layer itself, or a genuinely smaller theorem strictly inside the lift package.

\end{document}
